\documentclass[border=12pt]{standalone}
\usepackage[utf8]{inputenc}
\usepackage[T1]{fontenc}
\usepackage[american, siunitx]{circuitikz}

\sisetup{output-decimal-marker={,}}

\ctikzset{
    resistors/scale=0.8
}

\begin{document}
\begin{circuitikz}[font=\sffamily\small]

    % =========================================================================
    % 1. BLOCO ESQUERDO: ARDUINO UNO R3 - MESTRE (U1)
    % =========================================================================
    \draw[thick, fill=gray!5] (-6.0, -3.0) rectangle (-1.2, 3.0);
    \node[align=center, font=\sffamily\bfseries] at (-4.2, 0.4) {U1\\Arduino\\Uno R3};
    \node[font=\sffamily\scriptsize, text=gray!70!black] at (-4.2, -0.7) {(Mestre)};

    % Alimentacao de U1
    \draw (-3.6, 3.0) -- (-3.6, 3.8) node[vcc]{+5V};
    \node[below, font=\ttfamily\footnotesize] at (-3.6, 3.0) {5V};

    \draw (-3.6, -3.0) -- (-3.6, -3.8) node[ground]{};
    \node[above, font=\ttfamily\footnotesize] at (-3.6, -3.0) {GND};

    % Pinos I2C de U1
    \node[left, font=\ttfamily\scriptsize] at (-1.2, 0.8) {A4 (SDA)};
    \node[left, font=\ttfamily\scriptsize] at (-1.2, -0.8) {A5 (SCL)};

    % =========================================================================
    % 2. BARRAMENTO I2C (SDA, SCL) E RESISTORES DE PULL-UP (R1, R2 = 10k)
    % =========================================================================
    % Linhas SDA (bidirecional) e SCL (unidirecional)
    \draw[<->, >=stealth, thick] (-1.2, 0.8) -- (4.6, 0.8);
    \node[above, font=\ttfamily\scriptsize] at (-0.3, 0.8) {SDA};

    \draw[->, >=stealth, thick] (-1.2, -0.8) -- (4.6, -0.8);
    \node[below, font=\ttfamily\scriptsize] at (-0.3, -0.8) {SCL};

    % Barramento +5V dos resistores de pull-up (ponto apenas na juncao em T)
    \draw (0.6, 3.6) -- (2.8, 3.6);
    \node[vcc] at (1.7, 3.6) {+5V};
    \node[circ] at (1.7, 3.6) {};

    % Pull-up R1 (10 kOhm) na linha SDA
    \draw (0.6, 0.8) node[circ]{} -- (0.6, 1.4)
        to[R, l^=R1, a_=10<\kilo\ohm>] (0.6, 3.6);

    % Pull-up R2 (10 kOhm) na linha SCL
    \draw (2.8, -0.8) node[circ]{} -- (2.8, 1.4)
        to[R, l^=R2, a_=10<\kilo\ohm>] (2.8, 3.6);

    % =========================================================================
    % 3. BLOCO DIREITO: EEPROM 24LC256 - ESCRAVO (U2)
    % =========================================================================
    \draw[thick, fill=gray!5] (4.6, -3.0) rectangle (11.4, 3.0);
    \node[align=center, font=\sffamily\bfseries] at (8.3, 0.4) {U2\\24LC256};
    \node[font=\sffamily\scriptsize, text=gray!70!black] at (8.3, -0.5) {(EEPROM I2C)};

    % Alimentacao de U2 (VCC pino 8)
    \draw (8.0, 3.0) -- (8.0, 3.8) node[vcc]{+5V};
    \node[below, font=\ttfamily\scriptsize] at (8.0, 3.0) {VCC (8)};

    % Pinos I2C de U2
    \node[right, font=\ttfamily\scriptsize] at (4.6, 0.8) {SDA (5)};
    \node[right, font=\ttfamily\scriptsize] at (4.6, -0.8) {SCL (6)};

    % Pinos ligados ao GND (A0, A1, A2, VSS, WP)
    \node[above, font=\ttfamily\scriptsize] at (5.4, -3.0) {A0 (1)};
    \node[above, font=\ttfamily\scriptsize] at (6.7, -3.0) {A1 (2)};
    \node[above, font=\ttfamily\scriptsize] at (8.0, -3.0) {A2 (3)};
    \node[above, font=\ttfamily\scriptsize] at (9.3, -3.0) {VSS (4)};
    \node[above, font=\ttfamily\scriptsize] at (10.6, -3.0) {WP (7)};

    % Barramento comum de GND (sem pontos nos cantos externos de 90 graus)
    \draw (5.4, -3.0) -- (5.4, -4.0) -- (10.6, -4.0) -- (10.6, -3.0);
    \draw (6.7, -3.0) -- (6.7, -4.0) node[circ]{};
    \draw (8.0, -3.0) -- (8.0, -4.8) node[ground]{};
    \node[circ] at (8.0, -4.0) {};
    \draw (9.3, -3.0) -- (9.3, -4.0) node[circ]{};

\end{circuitikz}
\end{document}